Für die Ausarbeitung der Diplomarbeit stehen eine Vielzahl von Programmiersprachen zur Verfügung. Geachtet wurde bei der Auswahl auf die Einfachheit bzw. Vertrautheit der einzelnen Teammitglieder mit der Sprache, deren Funktionalität, das Bibliotheken-Angebot und die vorhandene Dokumentation. In dieser Diplomarbeit wird somit die Programmiersprache \gls{gls_python} verwendet. Auf die genauen Gründe zur Auswahl wird in dem folgenden Kapitel \ref{python_kapitel} eingegangen.

\subsection{\textit{Python}}\label{python_kapitel}
\begin{minipage}{0.6\textwidth}
	\gls{gls_python} ist eine vielseitig einsetzbare Programmiersprache. Die Erstellung einer ersten Version wurde 1989 von Guido van Rossum als Hobbyprojekt gestartet, während er am Centrum Wiskunde \& Informatica Institut in den Niederlanden arbeitete. Diese erste Version erschien 1991. Seither wird \gls{gls_python} stetig weiterentwickelt. Besonders durch seine einfache Syntax und die vielen vorhandenen Bibliotheken konnte sich \gls{gls_python} immer stärker auch in betrieblichen Umgebungen etablieren. Im Gegensatz zu Sprachen wie \gls{cprogrammiersprache}, C\#, etc. verwendet \gls{gls_python} keinen Compiler, um das Coding vor der Ausführung des Programms in Maschinencode umzuwandeln, sondern einen Interpreter, der beim Start des Programms jede Zeile nacheinander überprüft und ausführt. Diese Praktik wurde erstmals mit der Programmiersprache Lisp eingeführt und wird unter anderen von Java, Ruby oder PHP benutzt. 
	\cite[vgl.][]{Python_Software_Foundation:o.J., Pramanick_gfg:2019, Ryte:2021}
\end{minipage}%
\hfill
\begin{minipage}{0.37\textwidth}
	\centering	
	\includegraphics[width=0.58\textwidth]{Bilder/Python_logo}
	\captionof{figure}[\gls{gls_python} Logo (Quelle: \url{https://en.wikipedia.org/wiki/File:Python-logo-notext.svg}, Zugriff am 18.02.2024)]{\gls{gls_python} Logo (Quelle:\\
		\url{https://en.wikipedia.org/wiki/File:Python-logo-notext.svg}) \label{fig:python_logo}}
\end{minipage}
\vspace{1ex}

Die folgenden Vor- und Nachteile beziehen sich auf \textcite[vgl.][]{Ceaseo:2020}.
\paragraph{Vorteile}
\begin{itemize}
	\item Nicht nur die Syntax ist übersichtlich, sondern auch die Bibliotheken-Verwaltung ist einfach. Mit dem pip-Paketmanager können mit einem \mintinline{console}{pip install} alle vorhandenen Bibliotheken installiert werden.
	\item \gls{gls_python} kann in den unterschiedlichsten Anwendungen zum Einsatz kommen. Es können damit unter anderem Web-, Mobile- und Backendanwendungen erstellt werden. Außerdem eignet sie sich auch als Skriptsprache für Probleme, bei denen es keiner komplexen Software bedarf. Es können objektorientierte Strukturen genauso wie funktionale verwendet werden. 
	\item Wie bereits erwähnt, bietet \gls{gls_python} sehr viele unterschiedliche Bibliotheken, die von der Community ständig weiterentwickelt und verbessert werden. Außerdem besitzen diese Bibliotheken meist gut angelegtes Dokumentationsmaterial und Diskussionen bzw. Hilfestellungen in diversen Foren.
\end{itemize}

\paragraph{Nachteile}
\begin{itemize}{}{}
	\item Da die Sprache Zeile für Zeile interpretiert wird, dauert die Ausführung länger als bei kompilierten Sprachen.
	\item Zur Initialisierung einer Variable muss kein Datentyp angegeben werden. Dadurch hat der Programmierer weniger Schreibaufwand, allerdings können in erhöhtem Maße Laufzeitfehler auftreten und auch die Übersichtlichkeit ist betroffen. Genauso ist die Abwesenheit von Klammern eine kontroverse Funktionalität.
\end{itemize}
